\section{Preliminary Analysis}
\label{sec:preliminary}

We present three empirical observations that validate the physical causality gap and
motivate the design of PhysCausal.%
All measurements use a dedicated Android app that records \texttt{TYPE\_LINEAR\_ACCELERATION}
at \texttt{SENSOR\_DELAY\_FASTEST} while logging touch-event timestamps from the Android
input system.%
Human sessions are collected from 52 participants performing natural interaction tasks
(browsing, form filling, tapping prompted UI elements) while holding the phone in their
preferred grip.%
Agent sessions are generated by AutoGLM~\cite{autoglm2024} and AppAgent~\cite{appagent2024}
controlling identical UI workflows via ADB, with the phone resting on a flat surface.%

\subsection{Observation 1: Real Touch Produces a Detectable Accelerometer Impulse}
\label{sec:obs1}

Fig.~\ref{fig:obs1} shows the event-locked average of the linear acceleration magnitude
$|\delta(t)|$ aligned to touch-event timestamps ($t{=}0$) across 847 touch
events per participant (mean).%
For user-operated phones, a clear impulse transient of amplitude
$\approx$3.2\,mg appears immediately at $t{=}0$ and decays within
$\sim$60\,ms, consistent with the mechanical impulse response of finger contact transmitted
through the phone chassis~\cite{biomotouch2026}.%
The transient is present across all 52 participants and all
18 tested device models, confirming that physical causality is a
device-agnostic property of human touch.%
The signal-to-noise ratio of the transient above the ambient holding-motion background
is $\approx$8.4\,dB.%

For agent sessions (ADB-injected touch events), the event-locked average is flat at the
sensor noise floor throughout the $[-100,\,+200]$\,ms window.%
No impulse is present at any offset, confirming that software injection produces no
mechanical response.%

\begin{figure}[t]
  \centering
  \textcolor{green}{[Figure: Event-locked average of $|\delta(t)|$ around touch events.
  X-axis: time relative to touch event (ms), $-100$ to $+200$.
  Y-axis: linear acceleration magnitude (mg), 0 to 5.
  Two curves: Human (amber orange, clear impulse peak at $t{=}0$ reaching $\approx$3.2\,mg,
  decaying smoothly to baseline by $\sim$60\,ms) and LLM Agent (navy blue, flat at noise
  floor $\approx$0.2\,mg throughout).
  Shaded region $[0, 60]$\,ms in light amber. Error band ($\pm$1 std) around each curve.
  Vertical dashed line at $t{=}0$ labeled ``Touch Event.''
  White background, unified palette.]}
  \caption{Event-locked accelerometer response to touch events.
  Real finger contact (human) produces a detectable mechanical impulse ($\approx$3.2\,mg,
  SNR $\approx$8.4\,dB, decay $\lesssim$60\,ms);
  ADB-injected touch events (agent) produce no accelerometer response,
  confirming the structural physical causality gap.}
  \label{fig:obs1}
\end{figure}

\subsection{Observation 2: Touch Interval Distributions Are Separable}
\label{sec:obs2}

Fig.~\ref{fig:obs2} shows the empirical probability density of inter-touch intervals
$\Delta t_i = \tau_{i+1} - \tau_i$ for human and agent sessions.%
Human intervals follow a lognormal distribution
(mean 423\,ms, CV 0.68), consistent with neuromotor
variability in voluntary motor tasks~\cite{fitts1954}.%
AutoGLM and AppAgent sessions, which invoke a cloud vision-language model between each
action, exhibit a bimodal distribution: one peak at $\approx$130\,ms
corresponding to rapid repeat-tap actions, and a second peak at $\approx$1{,}050\,ms
corresponding to LLM inference plus network round-trip latency.%
ADB-script sessions without LLM inference show a near-uniform distribution with a
characteristic minimum inter-event interval floor of $\approx$47\,ms
imposed by the Android input dispatch scheduler.%

These distributional differences are complementary to the causality signal: an adversary
can adjust touch timing to mimic the human lognormal distribution (A2 adversary,
Section~\ref{sec:adversary}), but cannot simultaneously produce a physical impulse at
each artificially timed event.%

\begin{figure}[t]
  \centering
  \textcolor{green}{[Figure: Probability density of inter-touch intervals $\Delta t$.
  X-axis: interval duration (ms), log scale, 50 to 5000, ticks at 50, 100, 200, 500, 1000, 2000, 5000.
  Y-axis: probability density.
  Three curves: Human (amber orange, lognormal shape, single mode around 400--600\,ms),
  LLM Agent cloud (navy blue, bimodal: small peak at $\sim$130\,ms and larger peak at
  $\sim$1{,}050\,ms), ADB script (steel gray, near-uniform with hard floor at $\sim$47\,ms).
  Two vertical dashed navy lines marking the two LLM-agent peaks.
  White background, unified palette.]}
  \caption{Inter-touch interval distributions.
  Human intervals follow a lognormal distribution (mean 423\,ms, CV 0.68) from neuromotor
  variability; LLM agents show a bimodal signature from inference latency; scripted ADB
  shows a near-uniform distribution with a scheduler-imposed floor at $\approx$47\,ms.}
  \label{fig:obs2}
\end{figure}

\subsection{Observation 3: Joint Features Form Linearly Separable Clusters}
\label{sec:obs3}

To assess joint discriminability, we extract a two-dimensional feature vector per
session window: (i)~causality score $C = \frac{1}{N_e}\sum_{i=1}^{N_e} \mathbf{1}[|\delta(t_i^+)| > \theta]$,
the fraction of touch events followed by a detectable impulse, and
(ii)~bimodality index $B$ of the inter-touch interval distribution~\cite{bimodality1985}.%
Fig.~\ref{fig:obs3} shows the t-SNE projection~\cite{vandermaaten2008tsne} of
480 human windows and 320 agent windows across all
session types.%
The three classes (human, autonomous agent, overlay agent) form well-separated clusters
(silhouette score $\approx$0.81), confirming that even simple instantiations
of the two primitives yield strong separability.%
This motivates PCCN (Section~\ref{sec:pccn}), which learns richer representations
of both signals end-to-end.%

\begin{figure}[t]
  \centering
  \textcolor{green}{[Figure: t-SNE 2D scatter plot.
  Three clusters: Human (amber orange squares), Autonomous Agent (navy blue circles),
  Overlay Agent (steel gray triangles).
  Three visually separated cluster regions with minimal overlap.
  Human cluster: upper-left region.
  Autonomous Agent cluster: lower-right region.
  Overlay Agent cluster: upper-right region (partially overlapping with Human, but distinct).
  Axes: t-SNE Dim 1 / t-SNE Dim 2. No tick numbers (schematic).
  Faint dashed ellipse outline around each cluster in the corresponding color.
  Legend in top-right corner.
  White background, unified palette.]}
  \caption{t-SNE projection of joint causality score and bimodality index features.
  Human, autonomous agent, and overlay agent windows form three well-separated clusters
  (silhouette score 0.81), validating the discriminability of the two physical primitives
  across all threat classes.}
  \label{fig:obs3}
\end{figure}
